\documentclass[11pt,a4paper]{article}

\usepackage[utf8]{inputenc}
\usepackage[T1]{fontenc}
\usepackage[italian]{babel}
\usepackage[margin=2.2cm]{geometry}
\usepackage[hidelinks]{hyperref}
\usepackage{booktabs}
\usepackage{longtable}
\usepackage{tabularx}
\usepackage{array}
\usepackage{xcolor}
\usepackage{enumitem}
\usepackage{fancyhdr}
\usepackage{titlesec}
\usepackage{listings}

\definecolor{revalueGreen}{HTML}{2E7D32}
\definecolor{revalueBlue}{HTML}{1565C0}
\definecolor{revalueOrange}{HTML}{E65100}
\definecolor{lightGray}{HTML}{F5F5F5}
\definecolor{darkGray}{HTML}{333333}

\pagestyle{fancy}
\fancyhf{}
\lhead{\textbf{RE-VALUE} - Gruppo 21}
\chead{Milestone 4}
\rhead{7 giugno 2026}
\rfoot{\thepage}
\renewcommand{\headrulewidth}{0.4pt}

\titleformat{\section}{\Large\bfseries\color{revalueBlue}}{\thesection}{0.6em}{}
\titleformat{\subsection}{\large\bfseries\color{darkGray}}{\thesubsection}{0.6em}{}
\titleformat{\subsubsection}{\normalsize\bfseries\color{darkGray}}{\thesubsubsection}{0.6em}{}

\setlength{\parindent}{0pt}
\setlength{\parskip}{5pt}
\setlength{\headheight}{14pt}
\setlength{\emergencystretch}{2em}
\setlist[itemize]{leftmargin=1.3em,itemsep=2pt}
\setlist[enumerate]{leftmargin=1.5em,itemsep=2pt}
\renewcommand{\arraystretch}{1.18}

\newcolumntype{Y}{>{\raggedright\arraybackslash}X}
\newcolumntype{L}[1]{>{\raggedright\arraybackslash}p{#1}}
\newcommand{\todo}[1]{\textcolor{red}{\textbf{DA\_COMPLETARE: #1}}}
\newcommand{\ok}{\textcolor{revalueGreen}{\textbf{Completo}}}
\newcommand{\partiale}{\textcolor{revalueOrange}{\textbf{Parziale}}}
\newcommand{\proto}{\textcolor{revalueOrange}{\textbf{Prototipale}}}
\newcommand{\future}{\textcolor{darkGray}{\textbf{Futuro}}}
\newcommand{\code}[1]{\texttt{#1}}

\lstdefinestyle{plaincode}{
  basicstyle=\ttfamily\small,
  backgroundcolor=\color{lightGray},
  frame=single,
  breaklines=true,
  columns=fullflexible
}

\title{
  \vspace{-1.2cm}
  \textbf{\Huge RE-VALUE}\\
  \vspace{0.2cm}
  \Large Report Milestone 4 - Sprint 2
}
\author{
  Gruppo 21\\
  Alessandro Gremes \and Paolo Sarcletti \and Alessandro Turri
}
\date{2 giugno 2026}

\begin{document}

\maketitle
\thispagestyle{empty}

\begin{center}
\begin{tabular}{rl}
\textbf{Corso} & Ingegneria del Software - Universita' di Trento \\
\textbf{Anno accademico} & 2025/2026 \\
\textbf{Milestone} & 4 \\
\textbf{Deadline} & 07/06/2026 \\
\textbf{Repository} & \url{https://github.com/alegrem123/ReValue} \\
\textbf{API production host} & \url{https://revalue-backend-84jb.onrender.com} \\
\textbf{Frontend production} & \url{https://revalue-frontend.onrender.com} \\
\end{tabular}
\end{center}

\vfill

\begin{center}
\fbox{
\begin{minipage}{0.88\linewidth}
\textbf{Nota di consegna.} Il presente report segue il template ufficiale Milestone 4:
sezione introduttiva, sezione generale, sezione Sprint 2 e sezione finale. Le metriche
di test riportate sono quelle verificate localmente il 2 giugno 2026: backend 160/160,
frontend 4/4, mobile 4/4. I campi personali non disponibili nel repository sono marcati
come \code{DA\_COMPLETARE} e vanno compilati prima della consegna.
\end{minipage}
}
\end{center}

\newpage
\tableofcontents
\newpage

% ============================================================
\section{Sezione introduttiva}
% ============================================================

\subsection{Team members}

\begin{longtable}{L{3.4cm} L{3cm} L{4.2cm} L{4.2cm}}
\toprule
\textbf{Nome e cognome} & \textbf{Matricola} & \textbf{Account GitHub} & \textbf{Responsabilita' principali Sprint 2} \\
\midrule
Alessandro Gremes & \todo{matricola} & \todo{account GitHub} &
Sicurezza, CORS, rate limiting, notifiche, email service, dashboard admin, CI/CD. \\
\midrule
Paolo Sarcletti & \todo{matricola} & \todo{account GitHub} &
Recensioni, chat estesa, test di integrazione, QR edge cases, contratto API. \\
\midrule
Alessandro Turri & \todo{matricola} & \todo{account GitHub} &
Premi, segnalazioni, moderazione admin, deploy, documentazione e mobile refinements. \\
\bottomrule
\end{longtable}

\subsection{Project idea}

RE-VALUE e' una piattaforma web e mobile per favorire il riuso locale di oggetti.
Un utente pubblica un annuncio, un altro utente prenota l'oggetto, lo scambio
viene confermato tramite QR code e il sistema aggiorna automaticamente wallet,
storico, notifiche e reputazione. Il progetto integra anche premi riscattabili,
recensioni post-scambio, segnalazioni e strumenti di moderazione admin.

\subsection{Links}

\begin{longtable}{L{4cm} L{11cm}}
\toprule
\textbf{Risorsa} & \textbf{Link / riferimento} \\
\midrule
Repository GitHub & \url{https://github.com/alegrem123/ReValue} \\
\midrule
API Blueprint / Apiary & \todo{inserire link pubblico Apiary}. La fonte locale ufficiale e' \code{apiary.apib}. \\
\midrule
Backend deploy & \url{https://revalue-backend-84jb.onrender.com} \\
\midrule
Frontend deploy & \url{https://revalue-frontend.onrender.com} \\
\midrule
Endpoint smoke backend & \url{https://revalue-backend-84jb.onrender.com/api/v1/annunci} \\
\midrule
Documentazione OpenAPI secondaria & \code{docs/openapi.yaml}. Apiary resta la fonte di verita' richiesta dal corso. \\
\midrule
Collection Postman & \code{docs/postman/revalue-sprint2.postman\_collection.json} \\
\midrule
Account demo utente & \todo{email/password utente demo} \\
\midrule
Account demo admin & \todo{email/password admin demo} \\
\bottomrule
\end{longtable}

% ============================================================
\section{Sezione generale}
% ============================================================

\subsection{Strategia di branching}

La strategia Git del progetto evita una \emph{master only strategy}. Il repository
e' organizzato attorno a tre livelli:

\begin{itemize}
  \item \textbf{\code{main}}: branch stabile di consegna e deploy.
  \item \textbf{\code{preRelease}}: branch di integrazione, usato per aggregare le feature prima del merge finale.
  \item \textbf{\code{feat/<sigla>-<modulo>}}: branch individuali o di modulo, ad esempio feature di sicurezza, recensioni, premi, notifiche, mobile.
\end{itemize}

Il flusso adottato e' quindi:

\begin{lstlisting}[style=plaincode]
feat/<sigla>-<modulo>  ->  preRelease  ->  main
\end{lstlisting}

Rispetto allo Sprint 1, nello Sprint 2 la strategia e' stata rafforzata in tre modi:

\begin{itemize}
  \item le feature sono state collegate in modo piu' esplicito agli use case, ai requisiti e agli OCL;
  \item la CI e' stata usata come barriera di qualita' su test backend, frontend smoke e mobile smoke;
  \item gli endpoint legacy sono stati mantenuti ma deprecati formalmente, riducendo il rischio di rompere integrazioni gia' presenti.
\end{itemize}

\subsection{Product backlog}

Il backlog e' stato riorganizzato in modo da distinguere il flusso core, le feature
Sprint 2 e gli elementi progettuali/futuri. Questa separazione e' importante per
evitare che funzioni ridimensionate, come SSO o backend OpenStreetMap, vengano
presentate come completate.

\begin{longtable}{L{1.2cm} L{3.2cm} L{5.6cm} L{2.1cm} L{2.4cm}}
\toprule
\textbf{ID} & \textbf{Area} & \textbf{Descrizione} & \textbf{Priorita'} & \textbf{Stato} \\
\midrule
PB-01 & Auth & Registrazione, login, JWT, blocco utenti sospesi/bannati. & Alta & \ok \\
PB-02 & Annunci & CRUD annunci, filtri catalogo, stato, soft-delete. & Alta & \ok \\
PB-03 & Prenotazioni & Prenotazione oggetti disponibili, no self-booking, annullamento. & Alta & \ok \\
PB-04 & QR & Generazione e validazione QR per completamento scambio. & Alta & \ok \\
PB-05 & Wallet & Accredito/sottrazione crediti e storico transazioni. & Alta & \ok \\
PB-06 & Versioning API & Prefisso \code{/api/v1} per backend, frontend e mobile. & Alta & \ok \\
PB-07 & Recensioni & Recensioni post-scambio e reputazione pubblica. & Alta & \ok \\
PB-08 & Chat & Messaggi tra partecipanti, ricerca, polling incrementale, immagini opzionali. & Media & \ok \\
PB-09 & Notifiche in-app & Notifiche persistenti, non lette, mark-read singola/tutte. & Alta & \ok \\
PB-10 & Push native & Registrazione token Expo e invio via adapter backend. & Media & \partiale \\
PB-11 & Premi & Coupon, riscatti, stock illimitato con \code{stock=0}, codice univoco. & Media & \ok \\
PB-12 & Segnalazioni & Creazione segnalazione, no auto-segnalazione, motivo obbligatorio. & Alta & \ok \\
PB-13 & Moderazione admin & Dashboard, malus, sospensione, ban, riabilitazione, gestione annunci. & Alta & \ok \\
PB-14 & Supporto & Ticket supporto backend e API documentata. & Bassa & \partiale \\
PB-15 & Email & Adapter SMTP opzionale e fire-and-forget. & Bassa & \proto \\
PB-16 & SSO/SPID & Autenticazione esterna progettuale. & Bassa & \future \\
PB-17 & OpenStreetMap backend & Servizio backend dedicato per geocoding/mappa. & Bassa & \future \\
\bottomrule
\end{longtable}

\subsection{Definizione di ``Done''}

La Definition of Done adottata nello Sprint 1 e' stata mantenuta, ma nello Sprint 2
e' stata resa piu' rigorosa per includere API contract, test e tracciabilita':

\begin{enumerate}
  \item la feature e' implementata nel codice e raggiungibile tramite API REST versionata \code{/api/v1};
  \item gli endpoint sono documentati in \code{apiary.apib}, che e' la fonte ufficiale;
  \item eventuale OpenAPI secondaria e' coerente con Apiary per gli endpoint principali;
  \item sono presenti test unitari o di integrazione quando la feature ha logica backend;
  \item i vincoli OCL rilevanti sono applicati nel model, controller o service;
  \item frontend e mobile usano il client API centralizzato dove applicabile;
  \item la suite automatica passa senza regressioni;
  \item la documentazione finale distingue chiaramente tra completo, parziale, prototipale e futuro.
\end{enumerate}

% ============================================================
\section{Sezione Sprint 2}
% ============================================================

\subsection{Goal dello Sprint 2}

Il goal dello Sprint 2 era trasformare il prototipo Sprint 1 in un sistema piu'
maturo e difendibile: completare le feature sociali e di moderazione, stabilizzare
il contratto REST, migliorare sicurezza e deployment, introdurre notifiche, premi,
recensioni, segnalazioni e dashboard admin, mantenendo il flusso core
\emph{registrazione - annuncio - prenotazione - QR - wallet - recensione} sempre
funzionante e testato.

\subsection{Sprint 2 planning}

Lo Sprint 2 e' stato pianificato dal 18 maggio al 7 giugno 2026. Il piano operativo
ha assegnato le attivita' in modo coerente con le responsabilita' gia' assunte nel
progetto:

\begin{longtable}{L{2.1cm} L{4cm} L{6.2cm} L{2.2cm}}
\toprule
\textbf{Membro} & \textbf{Area} & \textbf{Task principali pianificati} & \textbf{Esito} \\
\midrule
AG & Sicurezza e notifiche &
CORS production whitelist, rate limiting, Helmet, Mongo sanitize, email service, notifiche backend, dashboard admin, CI. & \ok \\
\midrule
PS & Recensioni e chat &
Fix test regressione, OCL self-booking, recensioni, chat search/image/polling, QR edge cases, documentazione API. & \ok \\
\midrule
AT & Premi e moderazione &
Coupon, riscatti, segnalazioni, malus, admin backend, deploy, mobile premi/notifiche/recensione. & \ok \\
\midrule
Tutti & Integrazione &
Allineamento \code{/api/v1}, Apiary, test suite, D4, evidenze, deploy e demo. & \ok \\
\bottomrule
\end{longtable}

\subsubsection{Sprint 2 backlog}

\begin{longtable}{L{1.2cm} L{3.2cm} L{6.6cm} L{2cm} L{2cm}}
\toprule
\textbf{ID} & \textbf{Owner} & \textbf{Item} & \textbf{Priorita'} & \textbf{Stato} \\
\midrule
S2-01 & AG & Refactor endpoint a \code{/api/v1} su backend, web e mobile. & P0 & \ok \\
S2-02 & AG & CORS con branching \code{NODE\_ENV} e \code{FRONTEND\_URL}. & P0 & \ok \\
S2-03 & AG & Rate limiting auth con middleware dedicato, max 5, test 429. & P0 & \ok \\
S2-04 & AG & Email service Nodemailer opzionale, welcome email fire-and-forget. & P2 & \proto \\
S2-05 & AG & Notifiche backend: model, service, controller, route, mark-read. & P1 & \ok \\
S2-06 & AG & Hook notifiche su chat, prenotazioni, QR/scambio, segnalazioni. & P1 & \ok \\
S2-07 & AG & Helmet, Mongo sanitize e security tests. & P1 & \ok \\
S2-08 & PS & Fix catalog flow e test OCL self-booking. & P0 & \ok \\
S2-09 & PS & Recensioni backend e test integrazione. & P1 & \ok \\
S2-10 & PS & Chat estesa con ricerca, immagine, polling recente, typing. & P1 & \ok \\
S2-11 & PS & Apiary Sprint 1+2 allineato. & P0 & \ok \\
S2-12 & AT & Premi: coupon, riscatti, stock, OCL saldo. & P1 & \ok \\
S2-13 & AT & Segnalazioni e no auto-segnalazione. & P1 & \ok \\
S2-14 & AT & Moderazione admin: malus, sospensione, ban, riabilitazione. & P1 & \ok \\
S2-15 & AT & Deploy Render e documentazione deploy. & P1 & \ok \\
S2-16 & Tutti & Test cases, Postman, traceability matrix, D4 finale. & P0 & \ok \\
S2-17 & Tutti & Push native Expo end-to-end su dispositivo reale. & P2 & \partiale \\
S2-18 & Tutti & Load test reale 500 utenti con p95/throughput. & P2 & \partiale \\
\bottomrule
\end{longtable}

\subsubsection{Burndown chart}

Il burndown e' rappresentato con una stima a punti. I valori finali sono coerenti
con lo stato del repository al 2 giugno 2026: il core funzionale e' completo, mentre
rimangono come evidenze da completare il video demo, screenshot HD e benchmark reale
a 500 utenti.

\begin{center}
\begin{tabular}{lrrrrrrrr}
\toprule
\textbf{Data} & 18/05 & 21/05 & 24/05 & 27/05 & 30/05 & 02/06 & 05/06 & 07/06 \\
\midrule
Ideale & 72 & 60 & 48 & 36 & 24 & 12 & 6 & 0 \\
Reale / pianificato & 72 & 55 & 39 & 25 & 14 & 4 & 2 & 0 \\
\bottomrule
\end{tabular}
\end{center}

\begin{lstlisting}[style=plaincode]
Burndown reale (punti residui)
18/05 | ################################################ 72
21/05 | #####################################            55
24/05 | ##########################                       39
27/05 | #################                                25
30/05 | #########                                        14
02/06 | ###                                              4
07/06 | da chiudere: video demo, screenshot HD, benchmark reale
\end{lstlisting}

\subsection{Test cases: design e implementazione}

La verifica automatica e' uno degli elementi piu' solidi della Milestone 4. La run
locale del 2 giugno 2026 ha dato:

\begin{itemize}
  \item \textbf{Backend}: 17 test suite passate, 160 test passati, 0 fallimenti.
  \item \textbf{Frontend}: 4 smoke test passati.
  \item \textbf{Mobile}: 4 smoke test passati.
\end{itemize}

La matrice completa dei 74 test case Sprint 2 e' disponibile in
\code{docs/test-cases-sprint2.md}. La tabella seguente riporta i casi piu'
rappresentativi per area.

\begin{longtable}{L{1.7cm} L{2.6cm} L{6.2cm} L{3.2cm} L{2cm}}
\toprule
\textbf{ID} & \textbf{Modulo} & \textbf{Scenario} & \textbf{Fonte} & \textbf{Esito} \\
\midrule
TC-S2-001 & Auth & Registrazione valida con token JWT e utente senza password hash esposta. & \code{auth.test.js} & PASS \\
TC-S2-007 & Sicurezza & Route private rifiutano richieste senza token. & \code{security.test.js} & PASS \\
TC-S2-013 & Annunci & Creazione annuncio valido con scadenza futura e posizione. & \code{annunci.test.js} & PASS \\
TC-S2-028 & Prenotazioni & Donatore non puo' prenotare il proprio annuncio. & \code{prenotazioni.test.js} & PASS \\
TC-S2-030 & Concorrenza & Due acquirenti prenotano lo stesso annuncio: un successo e un conflitto. & \code{concurrent.test.js} & PASS \\
TC-S2-034 & QR & Acquirente valida QR e completa scambio con wallet aggiornati. & \code{swap.test.js} & PASS \\
TC-S2-049 & Premi & Riscatto premio con saldo sufficiente e codice univoco. & \code{premi.test.js} & PASS \\
TC-S2-056 & Segnalazioni & Auto-segnalazione rifiutata. & \code{segnalazioni.test.js} & PASS \\
TC-S2-059 & Recensioni & Recensione creata solo dopo scambio completato. & \code{reviews.test.js} & PASS \\
TC-S2-064 & Admin & Terzo malus produce auto-sospensione. & \code{admin\_moderation.test.js} & PASS \\
TC-S2-069 & Push token & Token Expo valido salvato su profilo utente. & \code{support\_push.test.js} & PASS \\
TC-S2-073 & Frontend & Smoke test viste principali e client API. & \code{frontend/tests} & PASS \\
TC-S2-074 & Mobile & Smoke test schermate, client API e wiring push token. & \code{mobile/tests} & PASS \\
\bottomrule
\end{longtable}

\subsection{Sprint review}

Durante la Sprint Review e' stato dimostrato il flusso end-to-end principale:

\begin{enumerate}
  \item registrazione e login utente;
  \item creazione annuncio da parte del donatore;
  \item visualizzazione nel catalogo e prenotazione da parte dell'acquirente;
  \item generazione QR da parte del donatore;
  \item validazione QR da parte dell'acquirente;
  \item aggiornamento wallet;
  \item recensione post-scambio;
  \item notifica e possibilita' di segnalazione;
  \item gestione admin di utenti, annunci e segnalazioni.
\end{enumerate}

Gli aspetti discussi dopo la demo sono stati:

\begin{itemize}
  \item il core applicativo e' stabile e verificato automaticamente;
  \item Apiary deve restare la fonte ufficiale del contratto API;
  \item SSO/SPID, backend OpenStreetMap ed email non devono essere presentati come feature complete;
  \item SHA-256 e' una scelta prototipale dichiarata e non una misura production-grade;
  \item push native e load test reale vanno documentati con prudenza, senza dichiarare misure non eseguite.
\end{itemize}

\subsection{Product backlog refinement}

Il refinement di fine Sprint 2 ha prodotto una distinzione netta tra:

\begin{itemize}
  \item \textbf{Core completo}: auth, annunci, prenotazioni, QR, wallet, recensioni, premi, segnalazioni, admin.
  \item \textbf{Estensioni complete come prototipo}: mobile Expo, notifiche in-app, deploy Render.
  \item \textbf{Parziali da dichiarare}: supporto UI, push su device reale, load test 500 utenti.
  \item \textbf{Futuro}: SSO/SPID, backend OpenStreetMap, hardening password con bcrypt/Argon2.
\end{itemize}

Il backlog futuro raccomandato e':

\begin{longtable}{L{1.2cm} L{5.5cm} L{7cm}}
\toprule
\textbf{ID} & \textbf{Item futuro} & \textbf{Motivazione} \\
\midrule
F-01 & Benchmark k6/autocannon a 500 utenti & Necessario per trasformare RNF2 da evidenza concorrente locale a metrica prestazionale. \\
F-02 & Screenshot HD e video demo & Evidenze richieste dalla consegna e utili per la commissione. \\
F-03 & UI supporto dedicata & Il backend supporto e' testato, ma manca una schermata completa dedicata. \\
F-04 & bcrypt/Argon2 con salt & Miglioramento production-grade rispetto a SHA-256 prototipale. \\
F-05 & SSO/SPID & Elemento progettuale da implementare solo in un contesto reale con provider esterno. \\
F-06 & Geocoding/backend OpenStreetMap & Riduzione del carico client e migliore validazione lato server. \\
\bottomrule
\end{longtable}

\subsection{Sprint retrospective}

\subsubsection{Cosa ha funzionato}

\begin{itemize}
  \item La separazione MVC/layered ha permesso di aggiungere moduli Sprint 2 senza riscrivere il core.
  \item I test di integrazione hanno protetto il flusso piu' rischioso: prenotazione concorrente e QR.
  \item Il prefisso \code{/api/v1} ha reso piu' chiaro il boundary REST.
  \item La documentazione Apiary e la traceability matrix hanno ridotto il rischio di incoerenza con D1/D2/D3.
  \item Il deploy Render ha reso disponibile una demo accessibile esternamente.
\end{itemize}

\subsubsection{Cosa non ha funzionato}

\begin{itemize}
  \item Alcune promesse progettuali iniziali erano troppo ambiziose per lo scope dello sprint.
  \item Alcuni task sono stati completati prima nel codice e solo dopo riallineati nei documenti.
  \item La forma delle risposte API non e' ancora perfettamente uniforme in tutti i controller.
  \item La misurazione prestazionale a 500 utenti non e' stata eseguita realmente entro la run del 2 giugno.
\end{itemize}

\subsubsection{Pratiche agile e dinamiche di gruppo}

Il gruppo ha applicato una forma pragmatica di Scrum/Kanban: backlog ordinato per priorita',
branch feature, integrazione progressiva su preRelease, review sui punti critici e retrospettiva
basata su evidenze. La divisione AG/PS/AT ha funzionato per ownership, ma l'integrazione finale
ha richiesto maggiore coordinamento su API contract, test e documentazione. La pratica piu'
utile e' stata trasformare le criticita' della review in task concreti e verificabili.

% ============================================================
\section{Sezione finale}
% ============================================================

\subsection{Diagramma del deploy finale}

\begin{lstlisting}[style=plaincode]
                         +-----------------------------+
                         |        Utente finale         |
                         |  Browser web / Expo mobile   |
                         +--------------+--------------+
                                        |
                                        v
              +-------------------------+-------------------------+
              |                                                   |
              v                                                   v
 +----------------------------+                   +----------------------------+
 | Render Static Site         |                   | Expo / Mobile client       |
 | Frontend HTML/CSS/JS       |                   | React Native app          |
 | https://revalue-frontend   |                   | API client /api/v1        |
 +-------------+--------------+                   +-------------+--------------+
               |                                                |
               +--------------------+---------------------------+
                                    |
                                    v
                    +---------------+----------------+
                    | Render Web Service             |
                    | Node.js + Express backend      |
                    | /api/v1 REST API               |
                    | JWT, CORS, Helmet, rate limit  |
                    +-------+-------------+----------+
                            |             |
                            v             v
              +-------------+--+       +--+------------------+
              | MongoDB Atlas |       | External services    |
              | Mongoose data |       | SMTP optional        |
              | Users, annunci|       | Expo Push optional   |
              | wallet, etc.  |       | OpenStreetMap client |
              +---------------+       +---------------------+
\end{lstlisting}

\subsection{Stack tecnologico utilizzato}

\begin{longtable}{L{3.2cm} L{4.6cm} L{7.2cm}}
\toprule
\textbf{Area} & \textbf{Tecnologie} & \textbf{Uso nel progetto} \\
\midrule
Backend runtime & Node.js, Express 5 & REST API, routing, middleware, server Render. \\
\midrule
Database & MongoDB Atlas, Mongoose & Persistenza documentale, schemi, validazioni, indici, transazioni dove necessario. \\
\midrule
Autenticazione & JWT, SHA-256 & Login/register, protezione route, digest password prototipale dichiarato. \\
\midrule
Sicurezza & CORS, Helmet, express-rate-limit, express-mongo-sanitize & Whitelist produzione, header HTTP, rate limit auth, mitigazione NoSQL injection. \\
\midrule
Email & Nodemailer & Adapter opzionale SMTP, fire-and-forget. \\
\midrule
Testing backend & Jest, Supertest, mongodb-memory-server & Unit/integration test con DB temporaneo. \\
\midrule
Frontend & HTML, CSS, JavaScript modulare, Bootstrap, Chart.js & Web app statica, dashboard admin, UI principali. \\
\midrule
Mobile & Expo SDK 54, React Native, AsyncStorage, Expo Camera, Expo Notifications, QRCode SVG & App mobile prototipale con schermate core, QR, notifiche e token storage. \\
\midrule
CI/CD & GitHub Actions, Render & Test automatici e deploy backend/frontend. \\
\midrule
Documentazione API & Apiary/API Blueprint, OpenAPI secondario, Postman collection & Contratto API ufficiale e supporto a demo/test manuali. \\
\bottomrule
\end{longtable}

\subsection{Conclusioni}

La Milestone 4 porta RE-VALUE da prototipo Sprint 1 a progetto completo e difendibile
per un esame di Ingegneria del Software. Il flusso core e' implementato, versionato,
testato e documentato: registrazione, pubblicazione annuncio, prenotazione, QR,
wallet, recensioni, notifiche, premi, segnalazioni e moderazione admin.

Il progetto dimostra coerenza architetturale con un modello MVC/layered: route,
controller, service, model e middleware hanno responsabilita' distinguibili. I test
automatici sono il principale elemento di maturita': 160 test backend, 4 smoke test
frontend e 4 smoke test mobile passano localmente.

I limiti residui sono stati dichiarati in modo esplicito: SHA-256 e' una scelta
prototipale, SSO/SPID e backend OpenStreetMap sono futuri, email e push dipendono
da servizi esterni, e il benchmark reale a 500 utenti deve essere prodotto prima
di dichiarare performance production-grade. Questa chiarezza rende il progetto
piu' credibile: non promette cio' che non e' stato misurato, ma mostra un core
solido e un piano realistico di evoluzione.

\subsection{Checklist finale prima della consegna}

\begin{longtable}{L{11cm} L{3cm}}
\toprule
\textbf{Attivita'} & \textbf{Stato} \\
\midrule
Compilare matricole e account GitHub dei membri. & \todo{} \\
Inserire link pubblico Apiary. & \todo{} \\
Inserire credenziali demo user/admin. & \todo{} \\
Registrarsi ad ESSE3 per la data dell'appello. & \todo{} \\
Compilare form individuale di autovalutazione per tutti i membri. & \todo{} \\
Consegnare video demo 2-3 minuti. & \todo{} \\
Allegare screenshot HD frontend/mobile se richiesti nel pacchetto finale. & \todo{} \\
Eseguire o dichiarare esplicitamente come non misurato il load test 500 utenti. & \partiale \\
Verificare ultima run test prima della consegna. & \ok \\
\bottomrule
\end{longtable}

\end{document}
